\section{Layered media: Love, Stoneley and Scholte waves}\label{sec:layered}
Rayleigh waves (\cref{sec:surface}) live on a single free surface. Two classical wave types need a \emph{second} medium: Love waves, guided by a surface layer over a half-space, and Stoneley waves, guided by the interface of two half-spaces. We give their isothermal elastic secular equations, check them against numbers printed in the original papers, and state how the heat law enters once thermal coupling is added. Everything here is the $\varepsilon\to0$ limit of the coupled problems catalogued in \cref{sec:atlas}; the coupled secular equations are not rederived.

\paragraph{Love waves.} Let a layer $0<z<T$ (shear speed $\beta_1$, rigidity $\mu_1$) rest on a half-space $z<0$ ($\beta_2$, $\mu_2$), with a traction-free upper surface. For anti-plane motion $u_y\propto e^{\mathrm i k(x-ct)}$ the layer carries $\cos$ and $\sin$ solutions with transverse number $kq_1$, $q_1=\sqrt{c^2/\beta_1^2-1}$, and the substrate carries a decaying exponential with $kq_2$, $q_2=\sqrt{1-c^2/\beta_2^2}$. Continuity of displacement and traction at $z=0$ and the free surface at $z=T$ give
\begin{equation}
\tan\!\big(kT\,q_1\big)=\frac{\mu_2}{\mu_1}\,\frac{q_2}{q_1},\qquad \beta_1<c<\beta_2 ,\label{eq:love}
\end{equation}
which is Love's relation \citep[eq.~(40), p.~162]{Love1911}. Love states that the root cannot be real unless the substrate shear speed exceeds that of the layer, so a Love wave requires $\beta_2>\beta_1$. Three properties follow from \eqref{eq:love}. (i) The $n$th mode appears at $c=\beta_2$ when $kT=n\pi/\sqrt{\beta_2^2/\beta_1^2-1}$ ($n=0,1,\dots$), so the number of modes at a given $kT$ is the integer part of $kT\sqrt{\beta_2^2/\beta_1^2-1}/\pi$ plus one. (ii) The fundamental mode has $c\to\beta_2$ as $kT\to0$ and $c\to\beta_1$ as $kT\to\infty$, decreasing monotonically. (iii) Dispersion is purely geometric: it is controlled by $kT$ and two ratios, with no intrinsic time scale. \Cref{fig:layered}(left) shows the first four modes for $\beta_2/\beta_1=\mu_2/\mu_1=2$. The script \path{check_layered.py} confirms (i)--(iii) numerically (mode counts $1,3,12$ at $kT=0.5,5,20$).

\paragraph{Stoneley waves.} Two half-spaces are welded at $z=0$; a wave $e^{\mathrm i k(x-ct)}$ decays into both. Each medium carries a dilatational and a shear potential decaying with $r_j=\sqrt{1-c^2/\alpha_j^2}$ and $s_j=\sqrt{1-c^2/\beta_j^2}$, so $c<\min(\beta_1,\beta_2)$ is necessary \citep[pp.~422--423]{Stoneley1924}. Continuity of the two displacement and two traction components gives a $4\times4$ determinant that must vanish \citep[eq.~(23), p.~418]{Stoneley1924}. Stoneley notes that for $\alpha_1=\alpha_2$ and $\beta_1=\beta_2$ a root exists at every density ratio, and that existence is not general when the wave speeds differ strongly. We solved the determinant independently, in real form with $k=1$. For incompressible media with $\rho_1=8.2$, $\rho_2=3.2$ and $\beta_1=\beta_2$ it gives $c=0.99287\,\beta$, the value printed by Stoneley (p.~422), and for $\rho_1=2\rho_2$, $\beta_2^2=2\beta_1^2$ it gives no real root, as printed (p.~423). For finite compressibility the root moves from $0.986\,\beta$ ($\alpha/\beta=1.5$) to $0.9929\,\beta$ ($\alpha/\beta\to\infty$). \Cref{fig:layered}(right) maps the existence region for incompressible media; it is a band around the line $\mu_1/\mu_2=\rho_1/\rho_2$ (equal shear speeds), widening toward large and small ratios.

\begin{figure}[t]\centering
\includegraphics[width=\textwidth]{figs/fig_layered.pdf}
\caption{Left: Love-wave phase speed against $kT$ for a layer on a stiffer half-space; modes enter at $c=\beta_2$. Right: existence of a real Stoneley root for two incompressible half-spaces as a function of density and rigidity ratios (dark: one root). Isothermal elastic baseline; computed with \texttt{fig\_layered.py}.}\label{fig:layered}
\end{figure}

\paragraph{Scholte's range of existence.} Scholte determined, for Stoneley waves, the region of the $(\rho_1/\rho_2,\mu_1/\mu_2)$ plane in which a wave exists, bounded by two curves; Stoneley waves ``can exist for every value'' of the ratios between them \citep[p.~120]{Scholte1947}. He treated a layer on a half-space with $T/L$ as parameter, $L$ being the transverse wavelength in the layer, and states that a surface-wave system exists for every density ratio and every $T/L$ provided the shear speed $\beta$ exceeds a Rayleigh speed (p.~123); the exact boundary curves of the one- and two-system regions are on p.~126. Our map in \cref{fig:layered} reproduces the shape of that region only for incompressible media; Scholte's curves are not recomputed here. For a liquid on a solid the same determinant, with the liquid carrying only a dilatational potential and zero shear traction, gives a Scholte-type interface wave slower than both the liquid sound speed and the solid Rayleigh speed. Water ($1480\,\mathrm{m/s}$) on steel ($\rho=7800\,\mathrm{kg/m^3}$, $\alpha=5960$, $\beta=3260\,\mathrm{m/s}$, handbook-type values) gives $c=0.9997$ times the water speed; the number is computed here and not taken from the source.

\paragraph{Anisotropy.} For bonded anisotropic half-spaces Barnett, Lothe, Gavazza and Musgrave write the existence condition as $\det\mathbf A=0$ at a subsonic speed $v_s$, where $\mathbf A$ is the sum of the two surface-impedance matrices, a $3\times3$ Hermitian matrix \citep[eq.~(2.11), p.~156]{Barnett1985}. This is the anisotropic form of the determinant condition above.

\paragraph{How the heat law enters.} In a thermoelastic layered problem each medium carries, besides the elastic partial waves, a thermal-type partial wave with $k_T^2=-\rho c\,s/G(s)$ in the template of \cref{sec:prelim}, and the interface conditions gain continuity of temperature and heat flux; the normal-incidence problem of \cref{sec:interfaces} has this structure. The heat law therefore enters the layered dispersion only through $G(s)$ in the thermal root, and the admissibility test $\operatorname{Re}G(\mathrm i\omega)\ge0$ of \cref{sec:admiss} applies to each medium separately. For piezoelectric media the analogous elastic baseline is the Bleustein--Gulyaev wave, with exact solutions for graded properties reported by \citet{Collet2006} (atlas entry A5). Coupled thermoelastic Love and Stoneley problems with specific heat laws are discussed in the sources of \cref{sec:atlas}; this review has not independently re-solved them.

\begin{table}[ht]\centering\small
\caption{Elastic baseline: values computed here from the secular equations against values printed in the original papers. The columns of Rayleigh's table are identified with the Poisson ratios by their values.}\label{tab:layered}
\begin{tabular}{@{}llcc@{}}\toprule
Quantity & Source & Printed & Computed here\\\midrule
Rayleigh speed, incompressible, $c/\beta$ & \citet{Rayleigh1885} & $0.9554$ & $0.9553$\\
Rayleigh speed, Poisson ratio $1/4$ & \citet{Rayleigh1885} & $0.9194$ & $0.9194$\\
Rayleigh speed, Poisson ratio $0$ & \citet{Rayleigh1885} & $0.8741$ & $0.8740$\\
Stoneley, $\rho_1=8.2$, $\rho_2=3.2$, $c/\beta$ & \citet{Stoneley1924} & $0.99287$ & $0.99287$\\
Stoneley, $\rho_1=2\rho_2$, $\beta_2^2=2\beta_1^2$ & \citet{Stoneley1924} & no real root & no real root\\\bottomrule
\end{tabular}
\end{table}
